\begin{longtable}[c]{|l|m{11cm}|}
    \caption{数据交换指令}
    \label{tab:pie-data-exchange-instructions}\\
    \hline
    \rowcolor{lightgray}
        \textbf{指令} & \textbf{定义} \\\hline
    \endfirsthead
    \rowcolor{lightgray}
    \textbf{指令} & \textbf{定义} \\\hline
    \endhead
    \endfoot
    \hyperref[instruction:MVQR]{MV.QR} & 将原 QR 寄存器中的值移动到目标 QR 寄存器中。
    \\\hline
    \hyperref[instruction:EEMOVI32A]{EE.MOVI.32.A} & 将 QR 寄存器中某片 32-bit 数据赋值给 AR 寄存器。
    \\\hline
    \hyperref[instruction:EEMOVI32Q]{EE.MOVI.32.Q} & 将 AR 寄存器内数据赋值给 QR 寄存器中某片 32-bit 数据空间。
    \\\hline
    \hyperref[instruction:EEVZIP8]{EE.VZIP.[8/16/32]} & 按照 1-字节/2-字节/4-字节 为处理单位对两个 QR 寄存器进行编码。
    \\\hline
    \hyperref[instruction:EEVUNZIP8]{EE.VUNZIP.[8/16/32]} & 按照 1-字节/2-字节/4-字节为处理单位对两个 QR 寄存器进行解码。
    \\\hline
    \hyperref[instruction:EEZEROQ]{EE.ZERO.Q} & 将指定 QR 寄存器清零。
    \\\hline
    \hyperref[instruction:EEZEROQACC]{EE.ZERO.QACC} & 将 QACC\_H 和 QACC\_L 寄存器清零。
    \\\hline
    \hyperref[instruction:EEZEROACCX]{EE.ZERO.ACCX} & 将指定 ACCX 寄存器清零。
    \\\hline
    \hyperref[instruction:EEMOVS8QACC]{EE.MOV.S8.QACC} & 按照 1-字节将 QR 寄存器分片，随后按符号位扩展成 20-bit 数据，赋值给 QACC\_H 和 QACC\_L 寄存器。
    \\\hline
    \hyperref[instruction:EEMOVS16QACC]{EE.MOV.S16.QACC} & 按照 2-字节将 QR 寄存器分片，随后按符号位扩展成 40-bit 数据，赋值给 QACC\_H 和 QACC\_L 寄存器。
    \\\hline
    \hyperref[instruction:EEMOVU8QACC]{EE.MOV.U8.QACC} & 按照 1-字节将 QR 寄存器分片，随后按 0 扩展成 20-bit 数据，赋值给 QACC\_H 和 QACC\_L 寄存器。
    \\\hline
    \hyperref[instruction:EEMOVU16QACC]{EE.MOV.U16.QACC} & 按照 2-字节将 QR 寄存器分片，随后按 0 扩展成 40-bit 数据，赋值给 QACC\_H 和 QACC\_L 寄存器。
    \\\hline
\end{longtable}